\documentclass[11pt]{article}
\addtolength{\oddsidemargin}{-1.cm}
\addtolength{\textwidth}{2cm}
\addtolength{\topmargin}{-2cm}
\addtolength{\textheight}{3.5cm}

\usepackage[pdftex]{graphicx}
\usepackage{pdflscape}
\usepackage{hyperref}
\usepackage{float}
\usepackage{cite}
\hypersetup{
	colorlinks=true,
	linkcolor=black,
	filecolor=magenta,
	urlcolor=cyan,
}

% define the title
\author{Team Longsword Notification}
\title{Phase 4 Testing Broadsword Notification}

\begin{document}
	\setlength{\parskip}{6pt}
	
	% generates the Cover page
	\begin{titlepage}
	
	\begin{center}
		% Upper part of the page       
		\includegraphics[width=0.7\linewidth]{Images/uniLogo.jpg}\\[1cm]    
		\textsc{\LARGE Longsword Notification}\\[0.3cm]
		% Title
		\rule{\linewidth}{0.5mm} \\[1cm]
		{ \huge \bfseries Phase 4 Testing Broadsword Notification}\\[0.5cm]
		\rule{\linewidth}{0.5mm} \\[1cm] 		
  
		\begin{minipage}{0.5\textwidth}
			\begin{flushleft} \large
				\emph{} \\
				Quinton {Swanepoel}
			\end{flushleft}
		\end{minipage}
		\begin{minipage}{0.4\textwidth}
			\begin{flushright} \large
				\emph{} \\
				15245510
			\end{flushright}
		\end{minipage}
		
		\begin{minipage}{0.5\textwidth}
			\begin{flushleft} \large
				\emph{} \\
				Frederick {Ehlers}
			\end{flushleft}
		\end{minipage}
		\begin{minipage}{0.4\textwidth}
			\begin{flushright} \large
				\emph{} \\
				11061112
			\end{flushright}
		\end{minipage}
		
		\begin{minipage}{0.5\textwidth}
			\begin{flushleft} \large
				\emph{} \\
				Tshepo {Malesela}
			\end{flushleft}
		\end{minipage}
		\begin{minipage}{0.4\textwidth}
			\begin{flushright} \large
				\emph{} \\
				14211582
			\end{flushright}
		\end{minipage}

		\begin{minipage}{0.5\textwidth}
			\begin{flushleft} \large
				Kyle {Erwin}
			\end{flushleft}
		\end{minipage}
		\begin{minipage}{0.4\textwidth}
			\begin{flushright} \large
				\emph{} \\
				15015302 
			\end{flushright}
		\end{minipage}
		
		\rule{\linewidth}{0.5mm} \\[1cm] 
		\textsc{\Large Stakeholders}\\[1cm]	
		
		\begin{minipage}{0.4\textwidth}
			\begin{flushleft} \large
				\emph{} \\
				Computer Science Department of University of Pretoria:
			\end{flushleft}
		\end{minipage}
		\begin{minipage}{0.4\textwidth}
			\begin{flushright} \large
				\emph{} \\
				Vreda Pieterse
			\end{flushright}
		\end{minipage}
		\begin{minipage}{0.4\textwidth}
			\begin{flushleft} \large
				\emph{} \\
				Team being tested:
			\end{flushleft}
		\end{minipage}
		\begin{minipage}{0.4\textwidth}
			\begin{flushright} \large
				\emph{} \\
				Broadsword Notification
			\end{flushright}
		\end{minipage}
		
	\end{center}
\end{titlepage}
	
	\tableofcontents
	
	\newpage
	
	\section{High level functional level requirements test models}
	\section{Lists of service contracts tested}
	\subsection{Pre-conditions}
	\begin{itemize}
		\item User details such as email address do exist and are valid.
		\item Notification request is complete with details of the email for example the subject and body of the email.
		\item The user receiving the email exists in the system
	\end{itemize}
	\subsubsection{Post-conditions}
	\begin{itemize}
		\item At least one notification has been sent to the appropriate user/users.
		\item System should log who Notifications are sent to, if the Notification request is not successful, the error message should be logged as well.
	\end{itemize}
	\textbf{The Notifications services should only fail under Three conditios that are:}
	\begin{itemize}
		\item Request has failed to meet the validation criteria specified in the post conditions.
		\item The request information is incomplete or has an error.
		\item There actually is no notification for the user.
	\end{itemize}
	\subsection{Team Broadsword Notification Serices}
	For this section we inspected the code of the Team Broadsword email services.
	\begin{itemize}
		\item In the code, the user details are supplied to send an email, this includes email, subject and body. 
		\item The email is not verified that it truelly is a real email address, therefore for fake user email, the code will produce exceptions. 
		\item The notification services of the Broadsword team does work, There are however no batch emails have been implemented.
		\item The system does have a method to log the email notifications, this is done by saving he notifications request object.  
	\end{itemize}
	\section{Evaluate the test cases of the previous team for non-functional requirements using your own criteria}	
		\subsection{Modularity}
			\paragraph{We felt that the broadsword notification team developed well written and modularized code. Each class has designed for a unique purpose and did nothing outside of that purpose. The functions themselves were not over complicated However they could have broken them up to into separate tasks.  \newline8/10}
		\subsection{In-code Comments}
			\paragraph{The comments for their implementation were decent however there were places were a comment could have been added or expanded upon. The classes themselves each have a descrption of what the classes are supposed to do but again were not exapned upon. \newline6/10}
		\subsection{Capacity}
			\paragraph{The team implemented the ability to send a single a email however they didn't provide the ability to send batch emails. \newline6/10}
		\subsection{Error Handling}
		    \paragraph{The error handling was implemented very well with in their python script, which is the dominant language they used. Every error thrown was caught and handled such as unsupported http requests. Their java script, however, did not provided any error handling on some object creations. This could lead to the function ending unexpectedly. \newline7/10}
		\subsection{Use Cases} 
			\paragraph{}

	\section{Add more tests if required}
		\paragraph{There were no tests found in the team Broadsword's Github repository. We would use the following python's built in libraries to test different cases}
		\subsection{Test cases for the NSQ server}	
			\begin{itemize}
				\item A test to ensure the NSQ server has succesfully started and is ready to be used
				\item A test to see if the program correctly handles it when the server failed to start
			\end{itemize}
		\subsection{Test cases for the notification server}
			\begin{itemize}
				\item A test to see that the Notification server has successfully started and is ready to be used
				\item A test to see if the program correctly handles it when the server failed to start	
			\end{itemize}
		\subsection{Test cases for the functions of the notification implementation}
			\begin{itemize}
				\item A test to see if a single email could be successfully be sent.
				\item A test to see if a batch of emails could be successfully be sent.
				\item A test to ensure the logging functionality successfully and in the correct order logged the requests
			\end{itemize}


	
	
\end{document}
